\section{The wild characteristic-three range}\label{sec:three}
This section proves Theorem~\ref{thm:three-range} for the wild defects and combines it with the tame criterion to obtain the largest possible characteristic-three level. Its valuation obstruction and wild invariant branch are proved in full below.

\subsection{Primitive torsion and its valuation sum}
We first isolate the valuation identity that applies before any descent condition is imposed.

\begin{lemma}\label{lem:valuation-sum}
Let $R$ be a complete discrete valuation ring with algebraically closed residue field of characteristic three. Suppose that a one-dimensional formal group of height two has full rational $3^n$-torsion over $\operatorname{Frac}(R)$. Write $d=v_R(3)$, and let $A$ and $B$ be the largest and smallest valuations of primitive $3^n$-torsion coordinates. Then $A\ge B\ge1$ and
\begin{equation}\label{eq:primitive-sum}
 d=2\,3^{n-1}\left(A+3^nB+
       \sum_{j=1}^{n-1}2\,3^{n-j-1}\min(A,3^{2j}B)\right).
\end{equation}
If $0\le r\le n-1$ and $3^{2r}B\le A\le3^{2r+2}B$, with no upper condition when $r=n-1$, then
\begin{equation}\label{eq:regional-sum}
 d=2\,3^{2n-r-2}(A+3^{2r+1}B).
\end{equation}
\end{lemma}

\begin{proof}
Choose primitive points $P,Q$ attaining $A,B$. When $A>B$, they form a basis over $\Z/3^n$: otherwise their reductions modulo three are dependent, so $Q=aP+3R_0$ for a scalar $a$ and a torsion point $R_0$. If $a$ is a unit, both summands have valuation greater than $B$; if it is not a unit, the same holds because multiplication by three strictly increases a positive formal valuation. Either alternative contradicts the choice of $Q$. When $A=B$, any basis has the required valuations.

We use Drinfeld's full-level ring with this marking. Its coordinate factorization, in \citet[Proposition~2.2.3(vii), pp.~439--441]{CarterStrauch2021}, has the following consequence:
\begin{equation}\label{eq:low-multiples}
 v_R(t(3^jQ))=3^{2j}B\qquad(0\le j<n).
\end{equation}
Indeed, after one level transition the coordinate of a primitive point is a unit times the product of its three primitive line lifts, each to the third power. Iterating $j$ times gives $3^j$ factors, each to exponent $3^j$. For the point $3^jQ$, these factors are represented by $aP+Q$ with $a\equiv0\pmod{3^{n-j}}$. Each has valuation $B$, since the contribution from $aP$ has strictly larger valuation. Specializing the universal identity proves~\eqref{eq:low-multiples} in $R$.

Primitive vectors modulo unit scalars form $\mathbf P^1(\Z/3^n)$. There are $3^n$ representatives $(a,1)$, all of valuation $B$, and the representative $(1,0)$ has valuation $A$. For each $1\le j<n$, there are $2\,3^{n-j-1}$ representatives $(1,b)$ with $v_3(b)=j$. Their valuations are $\min(A,3^{2j}B)$. If the two arguments of this minimum coincide, cancellation cannot raise the valuation: the resulting point is still primitive, whereas $A$ was maximal. The same observation handles $(a,1)$ when $A=B$.

Every line has $2\,3^{n-1}$ primitive generators, all with the same valuation. The power series $[3^n](T)/[3^{n-1}](T)$ has constant term three and its roots are precisely the primitive points. Weierstrass preparation therefore identifies $d$ with the sum of their valuations. This gives~\eqref{eq:primitive-sum}. Splitting its sum at $r$ and summing the two geometric series gives~\eqref{eq:regional-sum}. The formulas agree at a boundary $A=3^{2r+2}B$.
\end{proof}

\subsection{What a wild descent action forces}
We next explain why the largest primitive valuation has an extra factor of three when the descent group has wild order three.

\begin{lemma}\label{lem:additive-potential}
An elliptic curve over $K$ with full rational $3$-torsion and additive reduction is potentially good. After a finite good-reduction extension its special fibre is supersingular.
\end{lemma}

\begin{proof}
The Weil pairing gives $\mu_3\subset K$. If $v_K(j(E))<0$, Tate uniformization expresses $E$ as a quadratic twist of a split Tate curve. Its semisimplified three-torsion representation is
$\overline\epsilon\,\overline\chi_3\oplus\overline\epsilon$, where $\epsilon$ is the twisting character. Both $E[3]$ and $\overline\chi_3$ are trivial. Reduction modulo three is injective on $\{1,-1\}$, so $\epsilon=1$. The curve would have split multiplicative reduction, contrary to the hypothesis. Hence $j(E)$ is integral and $E$ is potentially good.

Over an algebraically closed residue field of characteristic three, the automorphism group of an ordinary elliptic curve is $\{1,-1\}$. Its nontrivial element acts nontrivially on the \emph{\'etale} three-torsion. A faithful nontrivial inertia action on an ordinary special fibre would therefore contradict full rational three-torsion. The special fibre is supersingular. We use here the good-reduction and specialization statements of \citet[Theorems~1--2 and Corollary~2 to Theorem~2, pp.~493, 496--497]{SerreTate1968}.
\end{proof}

\begin{lemma}\label{lem:wild-divisibility}
Suppose $E[3^n]$ is rational over $K$ and $E$ has semistability defect $m\in\{3,6,12\}$. Let $L/K$ be the minimal good-reduction extension after completed maximal unramified base change. In the good formal group over $\Ocal_L$, the largest primitive valuation $A$ is divisible by three.
\end{lemma}

\begin{proof}
First, the discriminant of every Weierstrass equation for $E$ is a cube in $K^\times$. To see this, use $E_*:y^2+y=x^3$, whose discriminant is $(-3)^3$ and whose three-torsion is rational over $K\supset\mu_3$. Since both three-torsion modules are constant symplectic planes, there is a symplectic isomorphism $E_*[3]\simeq E[3]$. The Hesse-pencil theorem puts $E$ in the pencil of $E_*$; the discriminant formula for that pencil is $\Delta(E_*)(\det A_t)^3$. These are exactly \citet[Theorem~1.1 and Proposition~5.1]{AnemaTopTuijp2018}. Changing a Weierstrass equation multiplies its discriminant by a twelfth power. Thus $\delta=v_K(\Delta)$ is divisible by three.

Let $\omega_K$ be a rational invariant differential and $\omega_L=u\omega_K$ an integral generator on the good model. The transformation formula for the discriminant gives
\[
 v_L(u)=\frac{m\delta}{12}.
\]
This is an integer divisible by three. Let $\sigma$ generate the wild subgroup and let $F=L^{\langle\sigma\rangle}$. Since the value group of $F^\times$ inside that of $L$ is $3\Z$, choose $b_0\in F^\times$ with $v_L(b_0)=v_L(u)$. Replacing $\omega_L$ by $(b_0/u)\omega_L$ gives a $\sigma$-invariant integral generator. Choose a formal coordinate with this differential at the origin, and make it odd by averaging under $[-1]$; two is a unit. The descent series has linear coefficient exactly one and only odd powers.

The supersingular curve is $y^2=x^3-x$. Its order-three automorphisms translate $x$ by a nonzero element of $\F_3$. Their first nonidentity term in an odd local coordinate has degree three. Hence the descent series is
\[
 t\longmapsto t+c t^3+O(t^5),\qquad c\in\Ocal_L^\times.
\]
For a rational nonzero torsion point with parameter $z$, this implies
\begin{equation}\label{eq:wild-point-displacement}
 v_L(\sigma z-z)=3v_L(z).
\end{equation}
The inverse descent convention changes $c$ by a unit and gives the same equality.

Let $b=v_L(\sigma\pi-\pi)-1$ be the positive lower break. Expanding $z$ in $\pi$ with coefficients in the fixed Cohen ring gives
\begin{equation}\label{eq:uniformizer-displacement}
 v_L(\sigma z-z)\ge v_L(z)+b,
\end{equation}
with equality if $3\nmid v_L(z)$. The lowest term gives equality in that case, and every later term changes at a strictly larger valuation. Applying the inequality to a point of minimum valuation $B$ gives $b\le2B$. If $A>B$ and $3\nmid A$, the equality at $A$ would give $b=2A>2B$, a contradiction.

It remains to consider $A=B$ and $3\nmid A$. Now $b=2A$. Write every primitive coordinate as $z=a_z\pi^A+O(\pi^{A+1})$. Comparing the leading terms of~\eqref{eq:wild-point-displacement} gives
\[
 \bar c\,a_z^3=d_0a_z,\qquad d_0\ne0,
\]
where $d_0=A\,\overline{(\sigma\pi-\pi)/\pi^{2A+1}}$ is independent of $z$. The roots of this additive polynomial form an $\F_3$-space of dimension at most one. On the other hand, for a primitive basis $P,Q$, the leading coefficient of $aP+bQ$ defines an injective $\F_3$-linear map from $\F_3^2$: a nonzero combination is primitive, so its leading coefficient cannot vanish without contradicting maximality of $A$. This contradiction proves $3\mid A$.
\end{proof}

\begin{proposition}\label{prop:necessary-wild}
Every curve in Theorem~\ref{thm:three-range} satisfies~\eqref{eq:wild-criterion}. The same numerical inequality holds for a tame defect $m=2$ or $4$.
\end{proposition}

\begin{proof}
Full rational $3^n$-torsion contains the cyclotomic extension, so $2\,3^{n-1}\mid e$ and $s\ge n-1$. The good model has $d=me$. Choose a region $r$ in Lemma~\ref{lem:valuation-sum}. For a wild defect, $A+3^{2r+1}B$ is divisible by three and $v_3(m)=1$. Equation~\eqref{eq:regional-sum} therefore implies
\[
 s\ge2n-r-2,\qquad r\ge\max(0,2n-2-s).
\]
For a tame defect, $v_3(m)=0$ and the integrality of the parenthesis gives exactly the same inequality. Since $A\ge3^{2r}B$ and $B\ge1$,
\[
 me=2\,3^{2n-r-2}(A+3^{2r+1}B)
       \ge8\,3^{2n+r-2}.
\]
Taking the smallest permitted $r$ proves the asserted condition.
\end{proof}

\subsection{A wild invariant curve in the level space}
We now construct an action-preserving one-dimensional specialization. The proof gives the complete local equation of the branch, rather than assuming that an invariant arc exists in a regular ring.

\begin{lemma}\label{lem:wild-arc}
Let $n\ge1$ and let $A\ge1$ be divisible by three. Set
\begin{equation}\label{eq:arc-degree}
 d=2\,3^{n-1}\left(A+3^n+
            \sum_{j=1}^{n-1}2\,3^{n-j-1}\min(A,3^{2j})\right).
\end{equation}
There is a complete discrete valuation ring $R$ with residue field $\overline{\F}_3$, $v_R(3)=d$, and a good supersingular elliptic curve over $R$ with full rational $3^n$-torsion and primitive valuations having extrema $A,1$.

The ring and marked curve admit a faithful semilinear $C_3$-action preserving every marked point. They admit such an action of $C_6$ if $A$ is odd, and of $C_3\rtimes C_4$ if $A\equiv1\pmod4$. The order-four generator acts by inversion on $C_3$. All actions are trivial on the residue field and induce the faithful indicated automorphisms of $y^2=x^3-x$.
\end{lemma}

\begin{proof}
Put $k=\overline{\F}_3$ and $W=W(k)$. A framed deformation of the supersingular curve has the form
\begin{equation}\label{eq:elliptic-family}
 \mathcal E_u:\quad Y^2=X^3+uX^2-X,\qquad u\in\mathfrak m.
\end{equation}
Indeed, completing the square is possible because two is invertible. The simple root lifting $X=0$ is unique; translate it to zero. A fourth root congruent to one then makes the coefficient of $X$ equal to $-1$. These operations are unique for a coordinate change reducing to the identity, leaving $u$ as the deformation parameter. The discriminant $16(u^2+4)$ is a unit. The Serre--Tate theorem identifies this deformation with the deformation of its height-two formal group.

Use the odd parameter $T=X^{-1/2}$ with sign chosen so that
\[
 Y=-T^{-3}(1+uT^2-T^4)^{1/2},\qquad
 \omega=\frac{dT}{(1+uT^2-T^4)^{1/2}}.
\]
Its logarithm begins $T-uT^3/6$, and comparison of $\log([3]T)$ with $3\log T$ gives
\begin{equation}\label{eq:cubic-coefficient}
 [3]_{T^3}=4u.
\end{equation}
This is an identity of integral coefficients, although it can be calculated after inverting three. On the special fibre, $[3](T)=-T^9$. One verification uses the three-division polynomial $3X^4+4uX^3-6X^2-1$: its tripling formula reduces to $X\mapsto X^9$, $Y\mapsto-Y^9$. For the latter sign, writing $f=X^3-X$, the characteristic-three identity
\[
 (f^2-1)(X-X^9)-f+f^5=0
\]
checks the addition formula directly.

Let $R_n$ be the universal full Drinfeld-level deformation ring. It is regular of dimension two, and the coordinates $x,y$ of a universal basis of $3^n$-torsion are regular parameters \citep[Theorem~2.1.2(ii), p.~437; Proposition~2.2.3(i), p.~439]{CarterStrauch2021}. The automorphisms of the special fibre act on the framing and hence on $R_n$, its elliptic curve, and its marking. Let $\sigma$ lift $X\mapsto X+1$, and let $\tau$ lift the automorphism for which
\[
 \tau u=-u,\qquad \tau T=iT,\qquad i^2=-1.
\]
These have orders three and four, with $\tau\sigma\tau^{-1}=\sigma^{-1}$. Their isomorphisms preserve the universal level structure. Modulo $(3,u)$ the series for $\sigma^j$, $j=0,1,2$, is
\[
 h_j(T)=\frac{T}{(1+jT^2)^{1/2}}.
\]
Because the origin is preserved, the difference of the two series lies in $(3,u)T\,W[[u,T]]$.

Set $M=3^{2n-2}$ and $H=6M$. The level-factorization theorem puts every three-torsion coordinate in $\mathfrak m_{R_n}^{M}$. In a Weierstrass factorization of $[3](T)$, its coefficient of $T^3$ contains at least six of the eight nonzero three-torsion coordinates in every term. Equation~\eqref{eq:cubic-coefficient} gives
\begin{equation}\label{eq:deformation-orders}
 u\in\mathfrak m_{R_n}^{H},\qquad
 3\in\mathfrak m_{R_n}^{8M}.
\end{equation}
In particular $H\ge6$. For $j<n$, the initial term of $[3^j]y$ is $(-1)^j y^{3^{2j}}$: the errors from $(3,u)$ have order at least $H$, whereas $3^{2j}\le M$.

Give $x,y,3$ weights $A,1,d$. The primitive coordinates have initial forms consisting of a nonzero scalar multiple of $Y$, or the lower-weight term of $X+c_jY^{3^{2j}}$, or this binomial when the weights coincide. These forms follow either from the level factorization used in Lemma~\ref{lem:valuation-sum}, or from the preceding exact special multiplication series and formal addition. At a weight equality the nonzero constants are in $\F_3$. Choose $c=i\in k$; then evaluation at $(X,Y)=(c,1)$ kills none of the initial factors because $i\ne\pm1$.

A Cohen presentation can be chosen as
\begin{equation}\label{eq:cohen-relation}
 R_n\simeq W[[X_0,Y_0]]/(3-P(X_0,Y_0)),
\end{equation}
where $P$ is a unit times the product of all primitive coordinates. The sum of their weights is $d$ by~\eqref{eq:arc-degree}. Its initial form has nonzero value $v_0$ at $(c,1)$. Each individual factor has weight less than $d$. Therefore a term in a lift of that factor which is divisible by three has larger weight than the factor, and cannot affect the weight-$d$ initial product. The coefficient of the linear term three in the relation is a unit; the relation generates the Cohen kernel.

Introduce
\[
 x=\pi^A Z,\qquad y=\pi,\qquad3=\pi^d V,
\]
and complete at $Z=c$, $V=v_0$, $\pi=0$. The chart has relations
\begin{equation}\label{eq:chart-relations}
 3=\pi^d V,\qquad V=P(\pi^A Z,\pi)/\pi^d.
\end{equation}
The quotient on the right is an integral convergent series. To verify this also for coefficients divisible by three, expand coefficients in Teichm\"uller digits and replace each $3^b$ by $\pi^{db}V^b$. All resulting exponents of $\pi$ are nonnegative. Its reduction is a polynomial in $Z$, independent of $V$, with value $v_0$ at $c$. Thus the second relation has derivative one in $V$ and eliminates $V$. The chart is a regular complete local ring of dimension two, with parameters $\pi,Z-c$. On it, $u\in\pi^H$ and $3\in\pi^d$, where $d\ge8M>H$.

The wild action extends to this chart before imposing any additional equation. The ratios $\sigma x/x$ and $\sigma y/y$ are units and are one modulo the maximal ideal. Hence
\[
 \sigma Z=(\sigma x/x)(\sigma y/y)^{-A}Z,\qquad
 \sigma V=(\sigma y/y)^{-d}V
\]
preserve its centre and its defining relations.

In the fraction ring of the chart, set $z_x=x^{-2}$, $z_y=y^{-2}$ and define the exact invariants
\begin{align}
 D_y&=\Tr_{\langle\sigma\rangle}(z_y^2),\notag\\
 I&=\frac{\Tr_{\langle\sigma\rangle}\bigl(z_y^2(z_x-z_y)\bigr)}{D_y},
 &U_y&=\Nm_{\langle\sigma\rangle}(y)^{-2}.
 \label{eq:trace-invariants}
\end{align}
The origin-preserving error bound gives
\[
 (\sigma^j x)^{-2}=x^{-2}+j+O(\pi^{H-2A}),\quad
 (\sigma^j y)^{-2}=y^{-2}+j+O(\pi^{H-2}).
\]
Consequently
\begin{align}
 D_y&=3y^{-4}+6y^{-2}+5+O(\pi^{H-4})
       =-1+O(\pi^{H-4}),\notag\\
 I&=z_x-z_y+O(\pi^{H-2A-4}),\notag\\
 U_y&=y^{-6}-y^{-2}+O(\pi^{H-6}).
 \label{eq:invariant-expansions}
\end{align}
For the estimate on $I$, first subtract $(z_x-z_y)D_y$ from its numerator; every remaining error has valuation at least $H-2A-4$. The norm estimate follows from
$\prod_{j=0}^2(w+j)=w^3+3w^2+2w$, whose difference from $w^3-w$ is divisible by three. Since $H>4$, $D_y$ is a unit of the chart. These are convergent integral expansions after the indicated normalizations.

Put $a=A/3$ and choose $\lambda\in W^\times$ reducing to $c^{-2}=-1$. Impose the invariant equation $I=\lambda U_y^a$. After multiplication by $x^2$ it is integral and becomes
\begin{equation}\label{eq:simple-branch}
 1-\pi^{2A-2}Z^2-\lambda Z^2(1-\pi^4)^a
                         +O(\pi^{H-4})=0.
\end{equation}
Its reduction is $1-\bar\lambda Z^2$ and has nonzero derivative $-2\bar\lambda c$ at $Z=c$. The complete local implicit-function theorem eliminates $Z$, as well as $V$ in~\eqref{eq:chart-relations}. The remaining ring is a complete DVR with uniformizer $\pi$ and $3=\pi^d$ times a unit. Every primitive coordinate has its assigned weight, so the extrema are $A,1$.

The equation defining the branch is exactly $C_3$-invariant, and the branch is unique over its residue point. It is therefore stable under $\sigma$. The wild action is faithful because
\[
 \sigma\pi=\pi-\tfrac12\pi^3+O(\pi^5).
\]
For the tame generator, conjugation permutes the trace and norm factors, giving $\tau I=-I$ and $\tau U_y=-U_y$. If $A\equiv1\pmod4$, then $a$ is odd, the equation is stable, and $\tau Z=i^{1-A}Z$ preserves the centre. Also $4\mid d$, so the centre for $V$ is stable. The action on $\pi$ is multiplication by $i$. For $\tau^2$, stability holds when $A$ is odd; then $2\mid d$. This gives the asserted groups and their faithful special actions. The universal marking is preserved at every step.
\end{proof}

\subsection{The integer criterion and the examples}
The invariant branch realizes all the integers permitted by the necessary bound.

\begin{proof}[Proof of Theorem~\ref{thm:three-range}]
Necessity is Proposition~\ref{prop:necessary-wild}. Assume~\eqref{eq:wild-criterion} for a chosen $m\in\{3,6,12\}$. Start with $r_0=\max(0,2n-2-s)$, and take the largest $r\in\{r_0,\ldots,n-1\}$ satisfying
\[
 me\ge8\,3^{2n+r-2}.
\]
Define
\begin{equation}\label{eq:construct-A}
 A=\frac{me}{2\,3^{2n-r-2}}-3^{2r+1},\qquad B=1.
\end{equation}
Since $v_3(me/2)=s+1$ and $2n-r-2\le s$, this is an integer divisible by three. The displayed inequality implies $A\ge3^{2r}$. If $r<n-1$, failure of the next inequality gives $A<3^{2r+2}$. Therefore~\eqref{eq:regional-sum} says precisely that the degree $d$ in Lemma~\ref{lem:wild-arc} is $me$. When $m=6$, the first term in~\eqref{eq:construct-A} is even, so $A$ is odd. When $m=12$, that term is divisible by four, and $3^{2r+1}\equiv3\pmod4$, so $A\equiv1\pmod4$.

Apply the lemma and take the fixed field of the indicated subgroup of order $m$. Its absolute index is $d/m=e$. The elliptic curve descends effectively, for example by using the invariant ample divisor $3[0]$. Its marked torsion points descend individually and are rational. Its faithful special inertia action has order $m$, so its exact semistability defect is $m$. It has additive reduction because it is potentially good and its inertia action is nontrivial.

Every additive elliptic curve in the question is covered by Lemma~\ref{lem:additive-potential}. Over $\overline{\F}_3$, the supersingular automorphism group is $C_3\rtimes C_4$: in $Y^2=X^3-X$, its elements have $X\mapsto a^2X+r$, $Y\mapsto a^3Y$, with $a^4=1$ and $r^3-r=0$. Thus the nontrivial possible defects are $2,3,4,6,12$. The tame possibilities satisfy the stronger bound in Proposition~\ref{prop:necessary-wild}; defect twelve gives the weakest bound and is always realized when that bound holds.

Writing $r=n-1$, the defect-twelve criterion is
\[
 r\le s,\qquad 2r\le L+1,\qquad 4r\le L+s+1.
\]
Its largest $r$ is the minimum in~\eqref{eq:maximum-three}. This proves the formula. A potentially good curve with full torsion above this level cannot be additive, and hence must be good. Existence at level $N(e)$ proves minimality.
\end{proof}

\begin{corollary}\label{cor:least-three-index}
The least absolute ramification index admitting an additive elliptic curve with full rational $3^n$-torsion, with arbitrary perfect residue field allowed, is
\[
 e_{\min}(n)=2\,3^{2n-2}.
\]
\end{corollary}

\begin{proof}
If $s\ge2n-2$, the divisibility $2\,3^s\mid e$ gives the bound. If $s\le2n-3$, the exponent $4n-4-s$ in the weakest, defect-twelve inequality is at least $2n-1$, and again $e\ge2\,3^{2n-2}$. At $e=2\,3^{2n-2}$, the defect-twelve criterion holds. Theorem~\ref{thm:three-range} gives an example.
\end{proof}

\begin{example}\label{ex:twenty-seven}
For $e=216$, one has $u=108$, $s=3$ and $L=4$, so $N(e)=3$. At $n=3$ and $m=12$, choose $r=1$ in~\eqref{eq:construct-A}; it gives $A=21$ and $B=1$. Direct substitution yields
\[
 d=18(21+27+6\cdot9+2\cdot21)=2592=12\cdot216.
\]
The invariant equation is $I=-U_y^7$ with $c=i$. The resulting degree-twelve descent gives full rational $27$-torsion with additive reduction. At $e=162$, the same calculation has $A=9$, $r=1$ and $d=1944$. This attains the least index in Corollary~\ref{cor:least-three-index}. For comparison, a tame defect at $e=216$ would require at least $8\,3^5/4=486$, so it cannot give full $27$-torsion there.
\end{example}
